% Generated draft: requires mathematical and source review.
\documentclass[11pt]{article}
\usepackage[margin=25mm]{geometry}
\usepackage{fontspec}
\setmainfont{Latin Modern Roman}
\usepackage{amsmath,amssymb,mathtools,unicode-math}
\setmathfont{Latin Modern Math}
\usepackage{enumitem,needspace,xurl,hyperref}
\setlength{\parindent}{0pt}
\setlength{\parskip}{5pt}
\setlength{\emergencystretch}{3em}
\newcommand{\N}{\mathbb{N}}
\newcommand{\Z}{\mathbb{Z}}
\newcommand{\Q}{\mathbb{Q}}
\newcommand{\R}{\mathbb{R}}
\newcommand{\C}{\mathbb{C}}
\newcommand{\F}{\mathbb{F}}
\newcommand{\A}{\mathbb{A}}
\newcommand{\PP}{\mathbb{P}}
\newcommand{\E}{\mathbb{E}}
\newcommand{\eps}{\varepsilon}
\newcommand{\dd}{\,\mathrm{d}}
\begin{document}
\section*{UnsolvedMath ID: 30002448}
\textbf{Source problem code:} catalogue entry
\subsection{Definitions and mathematical statement}
Let \(d,p\) be positive integers and let \(a<b\) be real endpoints, as implicit in the stated box domain. Set
\[
\Omega_d=(a,b)^d,\qquad \Delta=\sum_{j=1}^{d}\frac{\partial^2}{\partial x_j^2}.
\]
The source considers the formal eigenvalue problem
\[
-\Delta u(x)+V(x)u(x)=\lambda u(x)\quad (x\in\Omega_d),\qquad u|_{\partial\Omega_d}=0,\qquad u\not\equiv0.
\]
Here \(V\) is a potential whose admissible class is not specified. The boundary condition is homogeneous Dirichlet; the source does not specify the function space, operator domain, or classical or weak interpretation of the equation and boundary condition.

If an appropriate realization of this operator has a real spectrum bounded below with discrete eigenvalues, write its eigenvalues in nondecreasing order as \(\lambda_1\leq\lambda_2\leq\cdots\), conventionally counting multiplicities. This is conditional notation for the intended target \(\lambda_1,\ldots,\lambda_p\); the supplied record neither provides hypotheses ensuring these spectral properties nor specifies a multiplicity convention.

\textbf{Source question.} Can low-rank tensor methods overcome the curse of dimensionality in computing these \(p\) smallest eigenvalues, including for very large \(d\)?

\textbf{Unresolved qualifications.} This is a research objective rather than a fully quantified mathematical conjecture. The source does not fix a tensor representation or rank notion, an admissible family of potentials as \(d\) varies, a discretization, an accuracy criterion, or a computational cost model. Consequently, it does not specify what dependence on dimension, accuracy, \(p\), or the description of \(V\) would constitute overcoming the curse of dimensionality. No uniform complexity or tensor-rank guarantee is asserted by this draft. A precise theorem target requires these choices and the spectral hypotheses to be supplied.
\subsection{Short English statement}
The problem asks whether low-rank tensor methods can efficiently compute the \(p\) lowest eigenvalues of a Schr\"odinger-type operator on \((a,b)^d\) with homogeneous Dirichlet boundary conditions, even in very high dimension. The potential class, spectral assumptions, tensor format, accuracy requirement, and intended complexity bound remain unspecified, so the supplied question does not determine a precise yes-or-no conjecture.
\subsection{Sources}
\begin{itemize}
\item Local UnsolvedMath registry, supplied record with permanent ID \(30002448\), title ``Low-Rank Tensor Methods for High-Dimensional Eigenproblems,'' exact source problem code \texttt{catalogue entry}, in the collection ``Oberwolfach Reports Open Problems.'' No individual proposer is identified; the supplied proposed year is \(2013\).
\item Supplied report citation: \emph{Numerical Solution of PDE Eigenvalue Problems} (2014), Oberwolfach Reports, pp.\ 3221--3304, DOI \texttt{10.4171/owr/2013/56}. Supplied source URL and literature-source link, labeled ``original OWR report'': \url{https://doi.org/10.4171/owr/2013/56}. The background describes Oberwolfach Reports as documenting mathematical workshops and their open-problem discussions. No location within the report more specific than this page range is supplied.
\item Supplied literature assessment, checked on \texttt{2026-08-21}: the record is labeled ``open'' and characterized as a model and motivation for low-rank tensor eigensolvers rather than a precise theorem or yes-or-no claim; the assessment states that general guarantees therefore remain unresolved. This is supplied dated curation, not an independently verified literature assessment. The supplied background requests independent verification of time-sensitive status and cited claims.
\item The supplied registry labels the statement \texttt{corrected\_verified} and reports a manual semantic, typographic, and LaTeX review of the cleanup as self-contained and meaning-preserving. These are supplied editorial qualifications, not independent verification of the formulation, mathematical correctness, or current open status.
\end{itemize}
\end{document}
